\documentclass[11pt,a4paper]{article}
\usepackage[T1]{fontenc}
\usepackage{lmodern,microtype}
\usepackage[margin=25mm]{geometry}
\usepackage{amsmath,amssymb,amsthm,booktabs,array}
\usepackage[hidelinks]{hyperref}
\ifdefined\pdfinfoomitdate\pdfinfoomitdate=1\fi
\ifdefined\pdftrailerid\pdftrailerid{}\fi
\ifdefined\pdfsuppressptexinfo\pdfsuppressptexinfo=15\fi
\hypersetup{pdftitle={Sun's supercongruence for two Apery-like binomial sums},
 pdfauthor={Oleksiy Babanskyy},
 pdfsubject={A binomial supercongruence, finite differential invariants and modular periods},
 pdfkeywords={supercongruences, Apery-like sequences, Franel numbers, modular forms, complex multiplication}}
\pagestyle{plain}
\newtheorem{theorem}{Theorem}[section]
\newtheorem{proposition}[theorem]{Proposition}
\newtheorem{lemma}[theorem]{Lemma}
\newtheorem{corollary}[theorem]{Corollary}
\theoremstyle{remark}
\newtheorem{remark}[theorem]{Remark}
\newcommand{\Z}{\mathbb Z}\newcommand{\Q}{\mathbb Q}
\newcommand{\C}{\mathbb C}\newcommand{\Fp}{\mathbb F_p}
\newcommand{\HH}{\mathfrak H}\newcommand{\dd}{\partial}
\newcommand{\Tp}{\mathcal T_p}
\newcommand{\leg}[2]{\left(\frac{#1}{#2}\right)}
\newcommand{\smat}[4]{\begin{pmatrix}#1&#2\\#3&#4\end{pmatrix}}
\DeclareMathOperator{\ord}{ord}
\setlength{\emergencystretch}{2em}
\allowdisplaybreaks
\title{Sun's supercongruence for two Ap\'ery-like binomial sums}
\author{Oleksiy Babanskyy}
\date{}
\begin{document}
\maketitle
\begin{abstract}
Let $f_n=\sum_{j=0}^n\binom nj^3$ and
$Q_n=\sum_{k=0}^n\binom nk(-8)^{n-k}f_k$.
We prove that, for every odd prime $p$, the two sums
$\sum_{n=0}^{p-1}\binom{2n}n Q_n/(-32)^n$ and
$\sum_{n=0}^{p-1}\binom{2n}n Q_n/64^n$
are congruent modulo $p^2$ to $4x^2-2p$ when
$p=x^2+2y^2\equiv1,3\pmod8$, and to zero when
$p\equiv5,7\pmod8$.
For inert primes a finite quadratic differential invariant gives ordinary
and singular lifting results, including a formula for the defect at a
simple singularity. For split primes we certify the modular period by
identities of weight eight, construct an integral Hecke polynomial with
an explicit degree bound, and specialize a finite polynomial identity
at two complex multiplication points. The coefficient at degree $p$,
both classes of complex multiplication matrices, and the primes $3$ and
$5$ are treated explicitly.
\end{abstract}

\section{The supercongruence and the proof strategy}\label{sec:intro}
For $n\ge0$ define
\begin{equation}\label{eq:sequences}
 f_n=\sum_{j=0}^n\binom nj^3,\qquad
 Q_n=\sum_{k=0}^n\binom nk(-8)^{n-k}f_k,\qquad
 g_n=\binom{2n}n Q_n.
\end{equation}
The numbers $f_n$ are the Franel numbers. Write
\begin{equation}\label{eq:series-sums}
 G(t)=\sum_{n\ge0}g_nt^n,\qquad
 G_p(t)=\sum_{n=0}^{p-1}g_nt^n,\qquad
 S_-(p)=G_p(-1/32),\quad S_+(p)=G_p(1/64).
\end{equation}
The denominators are units for every odd prime. Congruences between
$p$-integral rational numbers are interpreted in $\Z_p$.
The subscript $p$ always denotes truncation at $p-1$.

\begin{theorem}\label{thm:main}
For every odd prime $p$,
\begin{equation}\label{eq:main}
 S_-(p)\equiv S_+(p)\equiv
 \begin{cases}
 4x^2-2p,&p=x^2+2y^2\equiv1,3\pmod8,\\
 0,&p\equiv5,7\pmod8
 \end{cases}
 \pmod{p^2}.
\end{equation}
\end{theorem}

This is Conjecture~3.6 of Zhi-Hong Sun
\cite[p.~2742]{sun}. The representation by $x^2+2y^2$ is unique up to
the signs of $x,y$, which do not affect $4x^2-2p$.
Its existence for $p\equiv1,3\pmod8$ follows from the norm-Euclidean
ring $\Z[\sqrt{-2}]$ and quadratic reciprocity.
Indeed, rounding both coordinates gives an error of norm at most
$1/4+2/4<1$, so this ring is Euclidean. For these primes a root of
$X^2+2$ modulo $p$ gives a proper nontrivial factor of $(p)$, and a
generator of that factor has norm $p$.

The sequence in \eqref{eq:sequences} satisfies the integral recurrence
\begin{equation}\label{eq:Q-recurrence}
 (n+1)^2Q_{n+1}=(-17n(n+1)-6)Q_n-72n^2Q_{n-1},
 \qquad Q_0=1,\quad Q_1=-6,
\end{equation}
as recorded in \cite[pp.~2729--2730]{sun}.
Sun's Theorem~3.7 on p.~2740 supplies the input
\begin{equation}\label{eq:Sun-input}
 S_-(p)\equiv S_+(p)\equiv0\pmod p
 \qquad(p>5,\ p\equiv5,7\pmod8).
\end{equation}
Although its printed statement includes $p>3$, the proof uses a sum
with denominator $100^n$. We use it only for $p>5$ and supply an
autonomous argument for $p=5$ in Section~\ref{sec:inert}.

Section~\ref{sec:invariant} proves the differential invariant and the
local lifting results. Section~\ref{sec:inert} completes the inert branch.
The endpoint congruence in Section~\ref{sec:endpoint} concerns individual
coefficients and is independent of the sums in \eqref{eq:main}.
Sections~\ref{sec:period} and~\ref{sec:geometry} certify the period,
its character, its divisors and the two complex multiplication values.
Section~\ref{sec:Hecke} proves the Hecke orbit and the integral polynomial
construction, following \cite[Lemma~3.7 and Section~5]{beukers}.
Section~\ref{sec:finite} derives the finite specialization, with the
degree-$p$ coefficient included. Section~\ref{sec:split} gives the
three required matrix constructions and completes the proof.
All uses of finite modular certificates are accompanied by the
holomorphy and valence argument that makes them sufficient.

\section{The finite differential invariant and local lifting}
\label{sec:invariant}
We first work with an integral sequence $u_n$, parameters $a,b,c\in\Z$,
and recurrence
\begin{equation}\label{eq:general-recurrence}
 u_0=1,\quad u_1=b,\qquad
 (n+1)^2u_{n+1}=(an(n+1)+b)u_n-cn^2u_{n-1}\quad(n\ge1).
\end{equation}
Let $v_n=\binom{2n}n u_n$, $V(t)=\sum_{n\ge0}v_nt^n$ and
$\theta=t\,d/dt$. The ratios of consecutive central binomial
coefficients give
\begin{equation}\label{eq:central-recurrence}
 n^3v_n=2(2n-1)(an(n-1)+b)v_{n-1}
 -4c(n-1)(2n-3)(2n-1)v_{n-2}\quad(n\ge1),
\end{equation}
where $v_{-1}=0$. Thus $LV=0$, with
\begin{equation}\label{eq:L}
 L=\theta^3-2t(2\theta+1)(a\theta^2+a\theta+b)
       +4ct^2(\theta+1)(2\theta+1)(2\theta+3).
\end{equation}
Put $D(t)=1-4at+16ct^2$. With coefficients on the left,
\begin{align}\label{eq:L-expanded}
 L={}&D\theta^3+(-6at+48ct^2)\theta^2\notag\\
 &+(-2(a+2b)t+44ct^2)\theta-2bt+12ct^2.
\end{align}
For a polynomial or formal series $Y$, define
\begin{equation}\label{eq:invariant}
 I(Y)=D\bigl(2Y\theta^2Y-(\theta Y)^2\bigr)
       +(\theta D)Y\theta Y+4t(-b+3ct)Y^2.
\end{equation}

\begin{lemma}\label{lem:Leibniz}
For every $Y$,
\begin{equation}\label{eq:Leibniz}
 \theta I(Y)=2YLY.
\end{equation}
In particular $I(V)=0$ in $\Q[[t]]$.
\end{lemma}
\begin{proof}
Set $E=4t(-b+3ct)$. The product rule gives
\[
 \theta I(Y)=2DY\theta^3Y+3(\theta D)Y\theta^2Y
 +(\theta^2D+2E)Y\theta Y+(\theta E)Y^2.
\]
The coefficients agree with $2YLY$ in \eqref{eq:L-expanded}.
Since $LV=0$, $I(V)$ is constant in characteristic zero, and its
constant term is zero.
\end{proof}

\begin{proposition}\label{prop:finite-invariant}
Under \eqref{eq:general-recurrence}, suppose all $u_n$ are integers.
For every prime $p>3$, the polynomial
$U(t)=\sum_{n=0}^{p-1}v_nt^n$ satisfies
\begin{equation}\label{eq:finite-invariant}
 I(U)\equiv0\pmod{p^2}\quad\text{in }\Z[t],
 \qquad \deg(U\bmod p)\le\frac{p-1}{2}.
\end{equation}
\end{proposition}
\begin{proof}
The exact residual of the truncation, from \eqref{eq:central-recurrence}, is
\begin{equation}\label{eq:LU}
 LU=-p^3v_pt^p+4cp(2p-1)(2p+1)v_{p-1}t^{p+1}.
\end{equation}
The omitted $v_p$ supplies the first term, and the highest remaining
degree is $p+1$. Since $p\mid\binom{2p-2}{p-1}$ and $u_{p-1}$ is
integral, $p\mid v_{p-1}$. Hence $LU\equiv0\pmod{p^2}$, and
\eqref{eq:Leibniz} gives $\theta I(U)\equiv0\pmod{p^2}$.

Here $\deg I(U)\le2p$ and $[t^0]I(U)=0$. For positive degrees other
than $p,2p$, division by the degree is division by a unit modulo $p^2$.
The two remaining coefficients must be computed before reduction.
As $V-U=v_pt^p+O(t^{p+1})$ and $I(V)=0$,
\begin{equation}\label{eq:resonance-p}
 [t^p]I(U)=-2p^2v_p.
\end{equation}
At this degree the only contribution of the omitted term pairs
$V(0)=1$ with $\theta^2(v_pt^p)$ in $2V\theta^2V$; all other terms
have an extra $t$ or a derivative of the constant.
At degree $2p$, only the highest coefficient of each factor $U$
contributes. Its scalar multiplier is
\[
 16c(p-1)^2+32c(p-1)+12c=4c(2p-1)(2p+1),
\]
so
\begin{equation}\label{eq:resonance-2p}
 [t^{2p}]I(U)=4c(2p-1)(2p+1)v_{p-1}^2.
\end{equation}
Both coefficients are divisible by $p^2$. This proves the polynomial
congruence in every degree.
Finally, $p\mid\binom{2n}n$ for $(p-1)/2<n<p$, proving the degree bound
after reduction modulo $p$.
\end{proof}

\begin{remark}\label{rem:upper-indices}
The half-degree bound in \eqref{eq:finite-invariant} holds modulo $p$.
The polynomial $U$ in the congruence modulo $p^2$ still includes every
index $0\le n<p$; its upper coefficients have not been discarded.
\end{remark}

\begin{lemma}\label{lem:multiplicities}
Let $\overline U\in\Fp[t]$ be the reduction of $U$ in
Proposition~\ref{prop:finite-invariant}, and let
$r\in\overline{\Fp}^{\times}$ be a zero of $\overline U$.
If $D(r)\ne0$, its multiplicity is exactly two.
If $D(r)=0$ and $D'(r)\ne0$, its multiplicity is exactly one.
\end{lemma}
\begin{proof}
At an ordinary zero the invariant modulo $p$ implies
$\theta\overline U(r)=0$, so the multiplicity is at least two.
The coefficients of the third and second ordinary derivatives in $L$ are
\begin{equation}\label{eq:ordinary-L-leading}
 A_3=t^3D,\qquad A_2=3t^2D+\tfrac32t^3D'.
\end{equation}
If the multiplicity were $3\le m<p$, write the leading local term as
$u(t-r)^m$, $u\ne0$. At degree $m-3$, the term $LU$ would have
coefficient $u r^3D(r)m(m-1)(m-2)\ne0$; the other derivatives
start in higher degrees. This contradicts $LU\equiv0\pmod p$.
The nonzero polynomial $\overline U$ has degree less than $p$, so
larger multiplicities cannot occur.

At a simple zero of $D$, suppose $m\ge2$. By the half-degree bound,
$m\le(p-1)/2$. In degree $m-2$, the combined leading coefficient
of the third and second derivatives is
\[
 u r^3D'(r)m(m-1)(m-\tfrac12).
\]
All remaining terms start at degree at least $m-1$.
This coefficient is nonzero: the bound excludes the resonance
$m=(p+1)/2$ as well as $m=0,1$. Hence the zero is simple.
\end{proof}

\begin{theorem}[Ordinary lifting]\label{thm:ordinary-lift}
In Proposition~\ref{prop:finite-invariant}, let $t_0\in\Z_p$ with
$t_0D(t_0)\in\Z_p^\times$. Then
\[
 U(t_0)\equiv0\pmod p\quad\Longrightarrow\quad
 U(t_0)\equiv0\pmod{p^2}.
\]
The conclusion holds for every $p$-integral lift of the same residue class.
\end{theorem}
\begin{proof}
Lemma~\ref{lem:multiplicities} shows that the reduced zero is exactly
double. Therefore $\theta U(t_0)\in p\Z_p$ and
$\theta^2U(t_0)\in\Z_p^\times$.
In \eqref{eq:finite-invariant}, all terms at $t_0$ except
$2D(t_0)U(t_0)\theta^2U(t_0)$ are divisible by $p^2$.
Its coefficient of $U(t_0)$ is a unit, proving the result.
The argument applies to every lift of the residue class.
\end{proof}

\begin{theorem}[Lifting at an exact simple singularity]
\label{thm:exact-singular-lift}
In Proposition~\ref{prop:finite-invariant}, let $t_0\in\Z_p^\times$
with $D(t_0)=0$ exactly and $D'(t_0)\in\Z_p^\times$. Then
$U(t_0)\equiv0\pmod p$ implies $U(t_0)\equiv0\pmod{p^2}$.
\end{theorem}
\begin{proof}
Lemma~\ref{lem:multiplicities} gives $U'(t_0)\in\Z_p^\times$.
Since $D(t_0)=0$, the invariant specializes to
\[
 U(t_0)\bigl[(\theta D)(t_0)\theta U(t_0)
              +4t_0(-b+3ct_0)U(t_0)\bigr]\equiv0\pmod{p^2}.
\]
The bracket is a unit: its first term is a unit and the second is
divisible by $p$. Division by this bracket proves the conclusion.
\end{proof}

\begin{theorem}[The defect at a simple singularity]\label{thm:defect}
In Proposition~\ref{prop:finite-invariant}, suppose $t_0\in\Z_p^\times$,
$D(t_0)\equiv U(t_0)\equiv0\pmod p$, and
$D'(t_0)\in\Z_p^\times$. Then $U'(t_0)$ is a unit and
\begin{equation}\label{eq:defect}
 U(t_0)\equiv D(t_0)\frac{U'(t_0)}{D'(t_0)}\pmod{p^2}.
\end{equation}
Consequently $U(t_0)\equiv0\pmod{p^2}$ if and only if
$D(t_0)\equiv0\pmod{p^2}$.
\end{theorem}
\begin{proof}
The unit assertion follows from Lemma~\ref{lem:multiplicities}.
The terms containing $D(t_0)U(t_0)$ or $U(t_0)^2$ vanish modulo $p^2$
in the invariant. Thus
\[
 -D(t_0)(\theta U(t_0))^2
 +(\theta D(t_0))U(t_0)\theta U(t_0)\equiv0\pmod{p^2}.
\]
Both $\theta U(t_0)=t_0U'(t_0)$ and
$\theta D(t_0)=t_0D'(t_0)$ are units. Dividing by them gives
\eqref{eq:defect}. The quotient $U'(t_0)/D'(t_0)$ is a unit, giving
the equivalence.
\end{proof}

\begin{corollary}\label{cor:root-discs}
For a zero $r\in\Fp^\times$ of $\overline U$, if $D(r)\ne0$, all
$p$ lifts modulo $p^2$ are zeros of $U$ modulo $p^2$.
If $D(r)=0$ and $D'(r)\ne0$, there is exactly one such lift; it is
the unique zero of $D$ modulo $p^2$ above $r$.
\end{corollary}
\begin{proof}
The ordinary assertion is Theorem~\ref{thm:ordinary-lift}.
In the singular case, Theorem~\ref{thm:defect} identifies the two sets
of zeros. For any lift $\widetilde r$, the congruence
\[
 D(\widetilde r+ps)\equiv D(\widetilde r)+psD'(\widetilde r)
 \pmod{p^2}
\]
has exactly one solution $s\in\Fp$.
\end{proof}

\section{The inert primes, including the prime five}\label{sec:inert}
For the recurrence \eqref{eq:Q-recurrence}, the parameters are
$(a,b,c)=(-17,-6,72)$, and
\begin{equation}\label{eq:D-Q}
 D(t)=1+68t+1152t^2=(1+32t)(1+36t).
\end{equation}
At $t_-=-1/32$, one has $D(t_-)=0$ exactly and
$t_-D'(t_-)=1/8$.
Theorem~\ref{thm:exact-singular-lift} applies for every $p>3$.
At $t_+=1/64$,
\begin{equation}\label{eq:D-plus}
 D(t_+)=75/32,\qquad D'(t_+)=104.
\end{equation}
For $p>5$ this is an ordinary point, so
Theorem~\ref{thm:ordinary-lift} applies.
The input \eqref{eq:Sun-input} therefore gives
$S_-(p)\equiv S_+(p)\equiv0\pmod{p^2}$ for all inert $p>5$.

For $p=5$, the first three original coefficients are
$Q_0=1,Q_1=-6,Q_2=42$. The half-degree bound gives, directly,
\begin{equation}\label{eq:p5-polynomial}
 G_5(t)\equiv1-12t+252t^2\equiv1+3t+2t^2\pmod5.
\end{equation}
This vanishes at $-1/32\equiv2$ and $1/64\equiv4\pmod5$.
The first point is handled by the exact singular lifting theorem.
For the second, \eqref{eq:D-plus} gives
$D(t_+)\equiv0\pmod{25}$ and $D'(t_+)\not\equiv0\pmod5$;
Theorem~\ref{thm:defect} gives $G_5(1/64)\equiv0\pmod{25}$.
Equivalently,
\[
 \frac1{64}-\left(-\frac1{36}\right)=\frac{25}{576}
\]
identifies this point modulo $25$ with the other exact singularity.
This calculation supplies the initial residue at five without using
the denominator-$100$ step in the proof of Sun's theorem.

\begin{remark}\label{rem:singular-shift}
Exact singular lifting need not persist under arbitrary shifts by $p$.
For example, direct evaluation at $p=7$, $t_-=-1/32$, gives
\[
 G_7(t_-)\equiv0\pmod{49},\qquad
 \theta G_7(t_-)\equiv5\pmod7,\qquad
 G_7(t_-+7)\equiv7\pmod{49}.
\]
The defect formula \eqref{eq:defect} describes the dependence on the lift.
\end{remark}

\section{The coefficient at degree \texorpdfstring{$p$}{p}}
\label{sec:endpoint}
The modular truncation will require $g_p\equiv g_1\pmod{p^2}$.
This is a congruence for individual coefficients, separate from
Theorem~\ref{thm:main}.

\begin{lemma}\label{lem:endpoint}
For every odd prime $p$,
\begin{equation}\label{eq:endpoint}
 Q_p\equiv-6\pmod{p^2},\qquad g_p\equiv-12=g_1\pmod{p^2}.
\end{equation}
\end{lemma}
\begin{proof}
The positive sequence F in Straub's notation is
\[
 A_F(n)=\sum_{k=0}^n(-1)^k8^{n-k}\binom nk f_k=(-1)^nQ_n.
\]
Straub \cite[p.~5, (10), Theorem~3.2]{straub} proves
$A_F(p^rn)\equiv A_F(p^{r-1}n)\pmod{p^{2r}}$ for
$p\ge3$ and $n,r\ge1$. Setting $n=r=1$ gives
$Q_p=-A_F(p)\equiv-A_F(1)=-6\pmod{p^2}$.
Also
\[
 \binom{2p}p=2\prod_{k=1}^{p-1}(1+p/k)\equiv2\pmod{p^2},
\]
because $\sum_{k=1}^{p-1}k^{-1}=0\pmod p$ for every odd prime.
Multiplication gives the second congruence.
\end{proof}

The required case can also be seen directly from the exact
representation of Gorodetsky \cite[p.~4, Proposition~1.2, (1.8)]{gorodetsky}:
\begin{equation}\label{eq:Gorodetsky}
 A_F(n)=\sum_{(\mathbf a,\mathbf b,\mathbf c,\mathbf d,\mathbf e)\in T(n)}
 (-1)^{a_1+b_2+c_3}
 \binom n{\mathbf a}\binom n{\mathbf b}\binom n{\mathbf c}
 \binom n{\mathbf d}\binom n{\mathbf e}.
\end{equation}
Here each of the five rows is a nonnegative triple of sum $n$, and
the three columns satisfy
\[
 a_i+b_i+c_i+d_i+2e_i=2n\qquad(i=1,2,3).
\]
At $n=p$, call a row pure if it is a permutation of $(p,0,0)$.
Every non-pure row has its multinomial divisible by $p$ and has a
coordinate nonzero modulo $p$. If it were the only non-pure row,
the corresponding column equation modulo $p$ would be impossible:
the pure rows contribute zero and the weight of this coordinate is
$1$ or $2$, both invertible. Therefore every term with a non-pure
row has at least two such rows and is divisible by $p^2$.

For the all-pure terms, division of every coordinate by $p$ gives
a bijection with $T(1)$. Their multinomials are $1$, and their signs
are preserved since $p$ is odd. Their sum is $A_F(1)=6$.
Thus \eqref{eq:Gorodetsky} proves the necessary endpoint as well.
Straub's proof uses this representation; these are two presentations
of the arithmetic step, with the representation itself attributed
to Gorodetsky. Neither presentation uses a value of $S_-(p)$ or $S_+(p)$.

\section{An exact modular period}\label{sec:period}
\subsection{Eta products, divisors and the cusp constants}
Let $q=e^{2\pi i\tau}$, $\eta_d=\eta(d\tau)$ and
$\dd=q\,d/dq=(2\pi i)^{-1}d/d\tau$. Define
\begin{gather}
 h=\frac{\eta_2\eta_6^5}{\eta_1^5\eta_3},\qquad
 z=\frac{\eta_1^6\eta_6}{\eta_2^3\eta_3^2},\qquad
 U_\eta=z^2,\qquad V_\eta=h^2z^2,\notag\\
 \Theta=U_\eta-72V_\eta=(1-72h^2)z^2,\qquad
 J=(\eta_1\eta_2\eta_3\eta_6)^2,\notag\\
 R=\frac{\eta_2^{12}\eta_3^{12}}{\eta_1^{12}\eta_6^{12}},\qquad
 H_2=\frac{\eta_2^{24}}{\eta_1^{24}}.\label{eq:eta-products}
\end{gather}
The subscript in $U_\eta$ distinguishes it from the finite polynomial $U$.

We use the eta-quotient criterion and cusp-order formula in
\cite[p.~2]{rousewebb}. For an exponent vector $(r_d)_{d\mid6}$,
the weight is $k=\frac12\sum r_d$. The congruences
\[
 \sum_{d\mid6}d r_d\equiv\sum_{d\mid6}(6/d)r_d\equiv0\pmod{24}
\]
and the condition that $\prod d^{r_d}$ is a rational square give
trivial character here, since every weight in the table is even.
For a cusp with denominator $c\mid6$, its local order is
\begin{equation}\label{eq:Ligozat}
 \frac{6}{24\gcd(c^2,6)}
 \sum_{d\mid6}\frac{\gcd(c,d)^2}{d}r_d.
\end{equation}
The cusp representatives are $0,1/2,1/3,\infty$, of widths
$6,3,2,1$, respectively. The resulting data are
\begin{center}\small
\begin{tabular}{c rrrr c c}
\toprule
 & $r_1$ & $r_2$ & $r_3$ & $r_6$ & weight & cusp orders\\
\midrule
$h$ & $-5$ & $1$ & $-1$ & $5$ & $0$ & $(-1,0,0,1)$\\
$R$ & $-12$ & $12$ & $12$ & $-12$ & $0$ & $(-1,1,1,-1)$\\
$H_2$ & $-24$ & $24$ & $0$ & $0$ & $0$ & $(-3,3,-1,1)$\\
$U_\eta$ & $12$ & $-6$ & $-4$ & $2$ & $2$ & $(2,0,0,0)$\\
$V_\eta$ & $2$ & $-4$ & $-6$ & $12$ & $2$ & $(0,0,0,2)$\\
$J$ & $2$ & $2$ & $2$ & $2$ & $4$ & $(1,1,1,1)$\\
\bottomrule
\end{tabular}
\end{center}
The nonnegative orders show that $U_\eta,V_\eta,J$ are holomorphic
modular forms. The other products are meromorphic modular functions.
Eta has no zeros on $\HH$. Thus $h$ has precisely one pole, simple,
on the compact curve $X_0(6)$, so it is a global coordinate on that
curve. This also establishes genus zero directly.

The exact phases, derived from the Dedekind transformation law in
Appendix~\ref{app:cusps}, give
\begin{equation}\label{eq:cusp-values}
 h(0)=\infty,\qquad h(1/2)=-1/9,\qquad
 h(1/3)=-1/8,\qquad h(\infty)=0.
\end{equation}
These agree with the data in \cite[Section~8.10]{beukers}.

\begin{lemma}\label{lem:eta-identities}
The following identities hold:
\begin{align}
 R&=h^{-1}+17+72h=\frac{(1+8h)(1+9h)}h,\label{eq:R}\\
 H_2&=\frac{h(1+9h)^3}{1+8h},\label{eq:H2}\\
 \Theta&=-\dd\log R.\label{eq:logR}
\end{align}
Moreover, $\Theta$ is holomorphic of weight two with trivial
character on $\Gamma_0(6)$, and is nonzero at every cusp.
\end{lemma}
\begin{proof}
In the coordinate $h$, $R$ has only simple poles at $0$ and $\infty$.
Consequently it has the form $a/h+b+ch$. The exact expansions
\[
 h=q+5q^2+19q^3+O(q^4),\qquad
 R=q^{-1}+12+78q+O(q^2)
\]
determine $a=1,b=17,c=72$. The three-dimensional rational function
space justifies the sufficiency of this comparison.

By \eqref{eq:cusp-values}, the right side of \eqref{eq:H2} has
the same divisor as $H_2$: its orders are $(-3,3,-1,1)$.
There are no other zeros or poles, and both functions have leading
term $q$ at infinity. Their quotient is a holomorphic nonvanishing
function on a compact curve, so it is the constant $1$.

The logarithmic derivative $\dd\log R$ is modular of weight two
on the base group. In the interior it is holomorphic since $R$
has no zeros or poles there. At a cusp of width $w$, the transformed
function has expansion $Q^r(c_0+O(Q))$, with
$Q=e^{2\pi i\tau/w}$, $r=\pm1$ and $c_0\ne0$.
Its logarithmic derivative has constant term $r/w$ and no negative
powers. Hence $\Theta+\dd\log R\in M_2(\Gamma_0(6))$.
The first three coefficients vanish, because both sides of
\eqref{eq:logR} begin $1-12q-12q^2$.

The index $[\mathrm{SL}_2(\Z):\Gamma_0(6)]=12$ gives valence bound
$2$ in weight two. A nonzero holomorphic form cannot have order
greater than $2$ at infinity; see
\cite[p.~2, Theorem~1.1(1)]{craig} for this classical bound.
The difference has order at least $3$ and is therefore zero.
At each cusp the constant term of $\Theta$ is $-r/w\ne0$.
This proves its nonvanishing without assuming that the two summands
$U_\eta$ and $72V_\eta$ cannot cancel.
\end{proof}

\subsection{Two identities of weight eight}
\begin{lemma}\label{lem:weight8}
Put $W=\Theta\dd\log J-2\dd\Theta$. Then
\begin{align}
 W^2-\Theta^4-68J\Theta^2-1152J^2&=0,\label{eq:P1}\\
 2\Theta\dd^2\Theta-3(\dd\Theta)^2+24J\Theta^2+864J^2&=0.
 \label{eq:P2}
\end{align}
\end{lemma}
\begin{proof}
For $M=\smat abcd\in\Gamma_0(6)$, set
$j=c\tau+d$ and $s=c/(2\pi i)$. The chain rule gives
\begin{align*}
 (\dd\Theta)(M\tau)&=j^4\dd\Theta+2sj^3\Theta,\\
 (\dd^2\Theta)(M\tau)&=j^6\dd^2\Theta
                   +6sj^5\dd\Theta+6s^2j^4\Theta,\\
 (\dd\log J)(M\tau)&=j^2\dd\log J+4sj.
\end{align*}
Therefore $W$ has weight four, and
$2\Theta\dd^2\Theta-3(\dd\Theta)^2$ has weight eight.
All derivative anomalies cancel exactly.

These covariant expressions are also holomorphic at the cusps.
For a cusp matrix $\sigma$, the same chain-rule computation
expresses their transforms in terms of the holomorphic expansions
of $\Theta|_2\sigma$ and $J|_4\sigma$. Derivatives introduce no
negative powers. The logarithmic derivative of $J|_4\sigma$ is
holomorphic there, since its expansion begins with a nonzero
constant times a power of the local parameter.
In the interior $J$ has no zeros. Thus both residuals really belong
to $M_8(\Gamma_0(6))$ before any coefficient comparison is made.

Appendix~\ref{app:certificates} gives the exact expansions and
an integer-product algorithm for the coefficients from $q^0$ to
$q^8$ of both residuals. All eighteen coefficients are zero.
The valence bound in weight eight is $8\cdot12/12=8$;
each residual has order at least $9$, so both vanish identically.
Only the classical holomorphic-form bound is used.
\end{proof}

\begin{proposition}\label{prop:period}
Let
\[
 P(h)=h(1+8h)(1+9h),\qquad K(h)=(1-72h^2)^2.
\]
Then
\begin{equation}\label{eq:parameter}
 t=\frac{J}{\Theta^2}=\frac{P(h)}{K(h)}
   \in q+q^2\Z[[q]],\qquad G(t(q))=\Theta(q).
\end{equation}
\end{proposition}
\begin{proof}
Comparison of eta exponents gives $J=h^2Rz^4$.
Equations~\eqref{eq:R} and \eqref{eq:eta-products} give the rational
expression for $t$. The products have integral expansions and
$\Theta(0)=1$, so $t\in q+q^2\Z[[q]]$.
This series has an integral formal inverse $q=q(t)$.

Let $e=\dd\log t=W/\Theta$, a unit formal series.
The first certificate implies
\[
 e^2=D(t)\Theta^2,\qquad D(t)=1+68t+1152t^2.
\]
Regard $Y(t)=\Theta(q(t))$ as a series in $t$.
Then $\dd=e\theta$, and logarithmic differentiation of the last
identity gives
\[
 2\frac{\theta e}{e}=\frac{\theta D}{D}
                       +2\frac{\theta Y}{Y}.
\]
Substituting $\dd=e\theta$ into \eqref{eq:P2}, dividing by $e^2$,
and using $J=tY^2$, gives
\[
 D\bigl(2Y\theta^2Y-(\theta Y)^2\bigr)
 +(\theta D)Y\theta Y+(24t+864t^2)Y^2=0.
\]
This is $I(Y)=0$ for $(a,b,c)=(-17,-6,72)$.
By \eqref{eq:Leibniz} and the fact that $Y$ is a unit,
$LY=0$ in characteristic zero.
The recurrence \eqref{eq:Q-recurrence} gives $LG=0$ for the same
operator. At each degree $n\ge1$, the equation determines the new
coefficient with multiplier $n^3\ne0$. The constant term $1$
therefore determines a unique formal solution, proving $Y=G$.
\end{proof}

The period also corresponds to the level-$6$ parametrization of
Chan--Cooper \cite{chancooper}. In their convention the recurrence
parameter is $c=-72$, the negative of the $c=72$ in
\eqref{eq:general-recurrence}. The proof above certifies the period
directly and does not require that parametrization as an input.

\section{Character, quotient geometry and complex multiplication values}
\label{sec:geometry}
For a matrix $M=\smat abcd$ of positive determinant, use
\begin{equation}\label{eq:slash}
 (Y|_2M)(\tau)=\det(M)(c\tau+d)^{-2}Y(M\tau).
\end{equation}
Scalar matrices act trivially in this convention.
Set
\begin{equation}\label{eq:AL-matrices}
 W_2=\smat{2}{-1}{6}{-2},\qquad
 W_6=\smat{0}{-1}{6}{0},\qquad
 \Gamma=\Gamma_0(6)+\langle2,6\rangle.
\end{equation}
We regard this full Atkin--Lehner extension projectively.
The normalizer action on $X_0(6)$ has a quotient of order four.
For completeness, the conjugation formulas in
Section~\ref{sec:Hecke} show that both $W_2,W_6$ normalize the base
group. Also $W_2^2=-2I$, $W_6^2=-6I$ and
\begin{equation}\label{eq:commutator}
 W_2W_6(W_6W_2)^{-1}=\smat{5}{2}{12}{5}\in\Gamma_0(6).
\end{equation}
Thus the projective quotient has at most four elements.
The determinant classes modulo squares in $\Q^\times$ are
$1,2,3,6$, so these four cosets are distinct.
This is the standard extension used in \cite[Section~3]{beukers}.

\begin{lemma}\label{lem:character}
The actions on $h$ are
\begin{equation}\label{eq:AL-actions}
 \phi_2(h)=-\frac{1+8h}{8(1+9h)},\qquad
 \phi_6(h)=\frac1{72h}.
\end{equation}
The form $\Theta$ has trivial character on $\Gamma_0(6)$ and
\[
 \Theta|_2W_2=-\Theta,\qquad \Theta|_2W_6=\Theta.
\]
The parameter $t$ is invariant under $\Gamma$.
\end{lemma}
\begin{proof}
Since $h$ is a global coordinate, an automorphism acts on it by a
M\"obius transformation. Direct action on the cusp fractions shows
that $W_2$ exchanges $\infty\leftrightarrow1/3$ and
$0\leftrightarrow1/2$, while $W_6$ exchanges
$\infty\leftrightarrow0$ and $1/2\leftrightarrow1/3$.
The four values in \eqref{eq:cusp-values} determine the two
transformations in \eqref{eq:AL-actions}; three values suffice in each case.
Exact rational substitution gives
\[
 R(\phi_2(h))=R(h)^{-1},\qquad R(\phi_6(h))=R(h),
 \qquad t(\phi_2(h))=t(\phi_6(h))=t(h).
\]
If $R(M\tau)=R(\tau)^\varepsilon$, the chain rule and
$d(M\tau)/d\tau=\det(M)/(c\tau+d)^2$ imply
\[
 (-\dd\log R)|_2M=\varepsilon(-\dd\log R).
\]
Equation~\eqref{eq:logR} proves the character assertions.
The signs follow from the transformation law of $R$.
\end{proof}

\begin{proposition}\label{prop:geometry}
The function $t$ is a Hauptmodul for $\Gamma$.
On $X_0(6)$ its two poles are double, and the zeros of $\Theta$
in the interior are simple and lie precisely above these poles.
The form $\Theta$ has no zeros at the cusps.
An invariant function with at most $m$ factors $\Theta^{-1}$ in
its denominator, holomorphic numerators and no other poles is a
polynomial in $t$ of degree at most $\lfloor m/2\rfloor$.
\end{proposition}
\begin{proof}
The polynomials $P$ and $K$ are coprime, and $P/K$ has degree four
as a map in the base coordinate $h$.
Its invariance under the quotient of order four makes its degree
on the quotient curve equal to one, proving the Hauptmodul assertion.
Its base poles are $h=\pm1/\sqrt{72}$, both double and both away
from the cusps. The formula $\Theta=(1-72h^2)z^2$, with $z$
nonzero in the interior, gives simple zeros precisely there.
The base group has no elliptic torsion: trace zero is impossible
modulo $3$, and trace $\pm1$ is impossible modulo $2$ for matrices
of determinant one with lower-left entry divisible by $6$.
Thus $h$ is an ordinary local coordinate at these points.
Nonvanishing at the cusps was proved in
Lemma~\ref{lem:eta-identities}.

The invariant function in the final assertion has no poles on the
quotient except at $t=\infty$, so it is a polynomial in $t$.
On the base curve its pole order is at most $m$.
A polynomial of degree $d$ in $t$ has pole order $2d$ there.
Hence $2d\le m$, as claimed. This calculation takes place on the
base curve and does not require integral local orders on the
orbifold quotient.
\end{proof}

\begin{lemma}\label{lem:CM-values}
For $\delta=\sqrt{-2}$ with positive imaginary part, set
\begin{equation}\label{eq:CM-points}
 \beta=\delta/2,\qquad\gamma=(2+\delta)/6.
\end{equation}
Then
\begin{equation}\label{eq:CM-values}
 t(\beta)=1/64,\qquad t(\gamma)=-1/32,
 \qquad \Theta(\beta)\Theta(\gamma)\ne0.
\end{equation}
\end{lemma}
\begin{proof}
Since $2\beta=-1/\beta$, the eta inversion formula
\cite[Section~23.18]{dlmf} gives
$H_2(\beta)=(-i\beta)^{12}=1/64$.
Equation~\eqref{eq:H2} therefore gives
\[
 64h(1+9h)^3-(1+8h)
 =(72h^2+16h+1)(648h^2+72h-1)=0.
\]
All eta factors are positive on the imaginary axis, so $h(\beta)>0$.
It follows that
$h(\beta)=(\sqrt6-2)/36$, the positive root of the second factor.
The identity
\[
 64P(h)-K(h)=-(8h^2-8h-1)(648h^2+72h-1)
\]
proves $t(\beta)=1/64$.

The equation $6\gamma^2-4\gamma+1=0$ shows that $W_2\gamma=\gamma$.
Its action on $h$ gives $72h(\gamma)^2+16h(\gamma)+1=0$.
Now
\[
 32P(h)+K(h)=(72h^2+16h+1)^2
\]
proves $t(\gamma)=-1/32$.
Neither of the two relevant quadratics has a common root with
$1-72h^2$. Thus the rational denominators and $\Theta$ are nonzero
at both points. These values follow from exact algebraic identities.
\end{proof}

\section{The Hecke orbit and an integral polynomial}\label{sec:Hecke}
\subsection{The orbit for a general prime}
The following argument makes explicit the level-$6$ case of
\cite[Lemma~3.7]{beukers}. It uses the weight-two transformation law
and the trivial character on the base group.
For $p>3$, let
\[
 \mathcal M(6,p)=\{M\in M_2(\Z):\det M=p,\ 6\mid M_{21}\}.
\]
\begin{lemma}\label{lem:orbit}
The left classes modulo $\Gamma_0(6)$ in $\mathcal M(6,p)$ have representatives
\begin{equation}\label{eq:Hecke-representatives}
 D_\infty=\smat p001,\qquad
 D_k=\smat1k0p\quad(0\le k<p).
\end{equation}
The group $\Gamma$ permutes the ratios
$r_D=(\Theta|_2D)/\Theta$ for these $p+1$ representatives.
\end{lemma}
\begin{proof}
For $M=\smat abcd\in\mathcal M(6,p)$, put $r=\gcd(a,c)>0$.
Then $r\mid p$, so $r=1$ or $p$, and $6\mid c/r$ because $p\nmid6$.
Choose $u,v$ with $ua+vc=r$. The matrix
\[
 L=\smat uv{-c/r}{a/r}\in\Gamma_0(6)
 \quad\text{gives}\quad
 LM=\smat r{b'}0{p/r}.
\]
A left translation reduces $b'$ modulo $p/r$.
This gives $D_\infty$ when $r=p$, and a $D_k$ when $r=1$.
Unimodular row operations preserve $r$. In the latter case,
$D_lD_k^{-1}=\smat1{(l-k)/p}01$ is integral if and only if
$k\equiv l\pmod p$, proving distinctness.

Right multiplication and conjugation by base-group elements preserve
$\mathcal M(6,p)$. For $M=\smat ab{6c}d$, direct multiplication gives
\begin{align}\label{eq:Hecke-conjugations}
 W_2^{-1}MW_2
 &=\begin{pmatrix}
 -2a-6b+6c+3d&a+2b-3c-d\\
 -6a-18b+12c+6d&3a+6b-6c-2d
 \end{pmatrix},\notag\\
 W_6^{-1}MW_6&=\smat d{-c}{-6b}a.
\end{align}
These matrices are integral with determinant $p$ and lower-left
entry divisible by $6$. The operations are reversible and hence
permute the left classes. The same formulas with determinant one
also verify normalization of $\Gamma_0(6)$ as used above.

For a normalizer element $w$, using the slash convention
\eqref{eq:slash} and its character $\chi(w)$ gives
\[
 r_D(w\tau)
 =\frac{\Theta|_2Dw}{\Theta|_2w}
 =\frac{\Theta|_2w(w^{-1}Dw)}{\chi(w)\Theta}
 =\frac{\Theta|_2(w^{-1}Dw)}{\Theta}.
\]
Reduction of $w^{-1}Dw$ to a representative $D_j$ uses a matrix
on the left in $\Gamma_0(6)$, whose character is trivial.
Thus these transformations permute the ratios.
The same computation for the base group completes the assertion
for the whole extension.
\end{proof}

\subsection{Polynomiality, degree and integral descent}
\begin{proposition}\label{prop:Hecke}
For every prime $p>3$ there is a polynomial
\[
 \Tp(X,t)=\sum_{m=0}^{p+1}C_m(t)X^m\in\Z[X,t]
\]
such that
\begin{equation}\label{eq:Hecke-product}
 \Tp(X,t(\tau))=
 \left(1-\frac{\Theta(p\tau)}{\Theta(\tau)}X\right)
 \prod_{k=0}^{p-1}
 \left(p^2-\frac{\Theta((\tau+k)/p)}{\Theta(\tau)}X\right).
\end{equation}
Its coefficients satisfy
\begin{equation}\label{eq:Hecke-properties}
 \deg C_m\le\lfloor m/2\rfloor,\qquad
 \Tp\equiv C_pX^p+C_{p+1}X^{p+1}\pmod{p^2},\qquad
 C_{p+1}\equiv1\pmod p.
\end{equation}
\end{proposition}
\begin{proof}
The product in \eqref{eq:Hecke-product} equals
\[
 p^{2p}\prod_{[D]}\left(1-\frac{\Theta|_2D}{p\Theta}X\right),
\]
since $r_{D_\infty}=p\Theta(p\tau)/\Theta(\tau)$ and
$r_{D_k}=\Theta((\tau+k)/p)/(p\Theta(\tau))$.
Lemma~\ref{lem:orbit} proves invariance of its coefficients under $\Gamma$.

The numerators $\Theta|_2D$ are holomorphic in the interior and
at the cusps. To justify the latter assertion, let
$\sigma\in\mathrm{SL}_2(\Z)$ represent a cusp, and choose
$\sigma'\in\mathrm{SL}_2(\Z)$ taking infinity to
the rational cusp $D\sigma\infty$.
Then $\sigma'^{-1}D\sigma=\smat ab0d$ with $a/d>0$.
Up to a nonzero constant, the expansion of
$(\Theta|_2D)|_2\sigma$ is that of $\Theta|_2\sigma'$
evaluated at $(a\tau+b)/d$.
Nonnegative exponents remain nonnegative; possible fractional
local exponents cause no poles.
The denominator $\Theta$ has nonzero constant term at every cusp.
Thus the coefficients have no cusp poles, and their only possible
interior poles lie above $t=\infty$.
Proposition~\ref{prop:geometry} makes them polynomials in $t$
and gives the degree bound $\lfloor m/2\rfloor$.

For integral descent, first work in
$\Z[\zeta_p][[s]]$, where $s=q^{1/p}$ and $\zeta_p$ is a primitive
$p$th root of unity. All coefficients of $\Theta$ are integers,
and its constant term is $1$, so its inverse has integral coefficients.
The symmetry $s\mapsto\zeta_ps$ of the product eliminates powers
of $s$ not divisible by $p$. Galois substitutions
$\zeta_p\mapsto\zeta_p^a$ permute the factors and make the resulting
coefficients rational. They are also algebraic integers, and hence
the coefficients lie in $\Z[[q]]$.
Since $t=q+q^2\Z[[q]]$, successive elimination of the lowest
coefficient of a polynomial in $t$ gives $C_m\in\Z[t]$.
More explicitly, the coefficient of $q^j$ after removal of the
terms of lower $t$-degree is exactly the coefficient of $t^j$,
and is an integer.

For $m<p$, every term contributing to $X^m$ contains at least
one constant factor $p^2$ from the $p$ finite factors.
After division by $p^2$, the same symmetries and integral descent
still apply. Thus $C_m\in p^2\Z[t]$ for $m<p$.

Since $p+1$ is even, the leading coefficient is
\begin{equation}\label{eq:leading-Hecke}
 C_{p+1}=
 \frac{\Theta(q^p)\prod_{k=0}^{p-1}\Theta(\zeta_p^ks)}
      {\Theta(q)^{p+1}}.
\end{equation}
The finite product in the numerator lies in $\Z[[q]]$.
Modulo $1-\zeta_p$, it equals $\Theta(s)^p$, which in
characteristic $p$ equals $\Theta(q)$.
On its already integral rational coefficients,
$(1-\zeta_p)\Z[\zeta_p]\cap\Z=p\Z$, so this is a congruence
modulo $p$ in $\Z[[q]]$.
Also $\Theta(q^p)\equiv\Theta(q)^p\pmod p$.
Equation~\eqref{eq:leading-Hecke} gives $C_{p+1}\equiv1$ as a series,
and triangular elimination gives the same congruence in $\Z[t]$.
\end{proof}

The ratio of a simple zero of $\Theta$ to a double pole of $t$
is $1/2$. The proof of Proposition~\ref{prop:Hecke} uses exactly
this bound and a modulus-$p^2$ product argument. It does not apply
the general modulus-$p^3$ theorem in \cite[Section~5]{beukers},
whose relevant stronger hypothesis is not satisfied here.

\section{The finite specialization and its degree boundary}\label{sec:finite}
\begin{theorem}\label{thm:finite-specialization}
For every prime $p>3$,
\begin{equation}\label{eq:finite-polynomial}
 C_p(t)+C_{p+1}(t)G_p(t)\equiv0\pmod{p^2}
 \quad\text{in }\Z_p[t].
\end{equation}
If $t_0\in\Z_p$ and $\alpha$ is an exact unit root in $\Z_p$
of $\Tp(X,t_0)$, then
\begin{equation}\label{eq:unit-specialization}
 G_p(t_0)\equiv\alpha\pmod{p^2}.
\end{equation}
\end{theorem}
\begin{proof}
The formal unit $\mathcal Q=\Theta(q)/\Theta(q^p)$ is a root of
\eqref{eq:Hecke-product}. Work in $\Z_p[[t]]$ using the integral
inverse $q=q(t)$.
Reducing by \eqref{eq:Hecke-properties} and dividing by
$\mathcal Q^p$ gives
\[
 C_p+C_{p+1}\mathcal Q\equiv0\pmod{p^2}.
\]
Write $C_{p+1}=1+pE$, with $E\in\Z[t]$.
Then
\begin{equation}\label{eq:comparison-polynomial}
 \mathcal Q\equiv-(1-pE)C_p\pmod{p^2},\qquad
 \deg\bigl((1-pE)C_p\bigr)
 \le\frac{p+1}{2}+\frac{p-1}{2}=p.
\end{equation}
The comparison polynomial can have degree $p$, as in Beukers'
boundary argument \cite[Theorem~6.1]{beukers} for ratio $1/3$
modulo $p^3$; here the ratio is $1/2$ modulo $p^2$.

Set $t^\sigma=t(q(t)^p)$.
The initial term $t(q)=q+O(q^2)$ implies
$t^\sigma=t^p+O(t^{p+1})$.
The exact period identity \eqref{eq:parameter} gives
\begin{equation}\label{eq:boundary-series}
 \mathcal Q=\frac{G(t)}{G(t^\sigma)}
 =G_p(t)+(g_p-g_1)t^p+O(t^{p+1}).
\end{equation}
Lemma~\ref{lem:endpoint} eliminates $g_p-g_1$ modulo $p^2$.
The polynomial in \eqref{eq:comparison-polynomial} and $G_p$
therefore have congruent coefficients at every degree from $0$
through $p$. Both have degree at most $p$, so they are congruent
as polynomials. Consequently $\mathcal Q\equiv G_p$ as a formal
series modulo $p^2$, which gives \eqref{eq:finite-polynomial}.

Now use only polynomials. The exact equation $\Tp(\alpha,t_0)=0$
and \eqref{eq:Hecke-properties} imply
\[
 \alpha^p\bigl(C_p(t_0)+C_{p+1}(t_0)\alpha\bigr)
 \equiv0\pmod{p^2}.
\]
Divide by the unit $\alpha^p$ and subtract
\eqref{eq:finite-polynomial} evaluated at $t_0$.
Since $C_{p+1}(t_0)\equiv1\pmod p$, it is a unit and
\eqref{eq:unit-specialization} follows.
No infinite series is evaluated at a $p$-adic unit parameter,
and no division by $G_p(t_0)$ is needed.
\end{proof}

\section{Complex multiplication matrices and the split primes}
\label{sec:split}
Let $p>3$ be split, and write $p=x^2+2y^2$ with $x,y>0$.
Then $x$ is odd, $\gcd(x,y)=1$, and exactly one of $x,y$ is
divisible by $3$. Indeed, if both were nonzero modulo $3$, then
$p\equiv1+2=0\pmod3$, and if both vanished then $9\mid p$.
Keep $\delta=\sqrt{-2}$ and the points \eqref{eq:CM-points}.

\subsection{The point \texorpdfstring{$\gamma$}{gamma}}
The integer $x+2y$ is coprime to $6y$: it is odd, is coprime to
$y$, and is nonzero modulo $3$.
Choose $A,B\in\Z$ satisfying $A(x+2y)+6By=1$, and put
\begin{equation}\label{eq:M-gamma}
 M_\gamma=\smat AB{-6y}{x+2y}\in\Gamma_0(6),
 \qquad k=B(x-2y)-Ay.
\end{equation}
The denominator of $M_\gamma\gamma$ is $x-y\delta$.
Rationalization gives
\begin{equation}\label{eq:gamma-image}
 M_\gamma\gamma=\frac{\gamma+k}{p}.
\end{equation}
To check both coordinates, the coefficient of $\delta$ in the
rationalized numerator is
$\{A(x+2y)+6By\}/6=1/6$.
The real part is $A(x-y)/3+Bx=1/3+k$, by the same B\'ezout identity.

The weight-two law has trivial character on the base group, so
\begin{equation}\label{eq:root-gamma}
 \alpha_\gamma=p^2\frac{\Theta(\gamma)}{\Theta(M_\gamma\gamma)}
 =\frac{p^2}{(x-y\delta)^2}=(x+y\delta)^2
\end{equation}
is an exact root of $\Tp(X,-1/32)$.
It comes from the factor indexed by $k\bmod p$ in
\eqref{eq:Hecke-product}; reduction of $k$ is valid since
$\Theta(\tau+1)=\Theta(\tau)$.

\subsection{The point \texorpdfstring{$\beta$}{beta} when
\texorpdfstring{$3\mid y$}{3 divides y}}
Choose $A,B\in\Z$ satisfying $Ax+2By=1$, and define
\begin{equation}\label{eq:M-beta-y}
 M_\beta=\smat AB{-2y}x\in\Gamma_0(6),\qquad k=Bx-Ay.
\end{equation}
The denominator is $x-y\delta$, and rationalization gives
$M_\beta\beta=(\beta+k)/p$.
For clarity, its real numerator is $Bx-Ay=k$, and its
$\delta$ coefficient is $Ax/2+By=1/2$.
Thus
\begin{equation}\label{eq:root-beta-y}
 \alpha_\beta=p^2\frac{\Theta(\beta)}{\Theta(M_\beta\beta)}
 =(x+y\delta)^2
\end{equation}
is an exact root of $\Tp(X,1/64)$.

\subsection{The point \texorpdfstring{$\beta$}{beta} when
\texorpdfstring{$3\mid x$}{3 divides x}}
Choose $A,B\in\Z$ satisfying $2Ay+Bx=1$, and define
\begin{equation}\label{eq:M-beta-x}
 M_\beta=\smat{2A}B{-2x}{2y},\qquad \det M_\beta=2,
 \qquad k=By-Ax.
\end{equation}
One has
\[
 M_\beta W_2^{-1}
 =\smat{-2A-3B}{A+B}{2x-6y}{-x+2y}\in\Gamma_0(6),
\]
where the lower-left entry is divisible by $6$ because $3\mid x$.
Therefore $\chi(M_\beta)=-1$.
The denominator is $2y-x\delta$, of norm $2p$.
Rationalizing $(A\delta+B)/(2y-x\delta)$ gives real part
$By-Ax=k$ and $\delta$ coefficient $(2Ay+Bx)/2=1/2$ after
division by $p$. Hence again $M_\beta\beta=(\beta+k)/p$.

The transformation convention \eqref{eq:slash} gives
\[
 \Theta(M\tau)=\frac{\chi(M)}{\det M}(c\tau+d)^2\Theta(\tau).
\]
In this case it yields the exact root
\begin{equation}\label{eq:root-beta-x}
 \alpha_\beta=p^2\frac{\Theta(\beta)}{\Theta(M_\beta\beta)}
 =-\frac{2p^2}{(2y-x\delta)^2}=(x-y\delta)^2
\end{equation}
of $\Tp(X,1/64)$.
The last equality uses $2y-x\delta=-\delta(x+y\delta)$ and
$\delta^2=-2$. Both the negative character and the determinant
factor are part of this computation.

\subsection{Separate embeddings and the rational trace}
Each root in \eqref{eq:root-gamma}, \eqref{eq:root-beta-y} and
\eqref{eq:root-beta-x} has been constructed exactly in $\Q(\delta)$.
Its polynomial equation has rational coefficients, so it remains
valid under either embedding into $\Q_p$.
Since $\leg{-2}p=1$, there are two such embeddings.
The two factors $x+y\delta$ and $x-y\delta$ have product $p$;
they cannot both be divisible by $p$, since their sum is $2x$
and $p\nmid x$. Thus one is a unit and the other has valuation one.
Choose the embedding in which the factor occurring in the required
root is a unit. Its conjugate square is zero modulo $p^2$, and
\begin{equation}\label{eq:CM-trace}
 (x+y\delta)^2+(x-y\delta)^2=4x^2-2p
\end{equation}
shows that the unit root is congruent to $4x^2-2p\pmod{p^2}$.

When $3\mid x$, the roots constructed at $\gamma$ and $\beta$
are conjugates. They are not both units in a single embedding.
For $S_-(p)$ choose the embedding making \eqref{eq:root-gamma}
a unit; for $S_+(p)$ choose the other embedding making
\eqref{eq:root-beta-x} a unit.
The sums are rational $p$-integral numbers, and both embeddings
give the same rational trace in \eqref{eq:CM-trace}.
Theorem~\ref{thm:finite-specialization}, applied separately to the
two polynomial equations, proves
\[
 S_-(p)\equiv S_+(p)\equiv4x^2-2p\pmod{p^2}
 \qquad(p>3,\ p\equiv1,3\pmod8).
\]
Unit roots have therefore been obtained from the matrices and the
norm before applying the congruence for the finite sums.

\subsection{The prime three and completion of the proof}
For $p=3$, the definitions give $Q_0=1,Q_1=-6,Q_2=42$.
Both sums have addends $1,6,0$ modulo $9$, and hence both equal
$7\pmod9$. This is $4-2\cdot3$ for $3=1^2+2\cdot1^2$.
As an additional direct check at five, the original addends are
\begin{center}
\begin{tabular}{c l c}
\toprule
sum & addends modulo $25$ & total\\
\midrule
$S_-(5)$ & $1,16,23,5,5$ & $0$\\
$S_+(5)$ & $1,17,12,15,5$ & $0$\\
\bottomrule
\end{tabular}
\end{center}
Section~\ref{sec:inert} proved the inert case for every $p>3$,
and the preceding matrix argument proved the split case for $p>3$.
The direct evaluation at three completes Theorem~\ref{thm:main}.

\section{Further consequences and the precision of the result}
\label{sec:consequences}
\subsection{The squarefree part of a finite period}
\begin{proposition}\label{prop:squarefree}
Under Proposition~\ref{prop:finite-invariant}, suppose $D\bmod p$
is separable. There are $S,V\in\Fp[t]$ such that
\[
 \overline U=S V^2,\qquad S(0)=V(0)=1,
 \qquad S\mid D,\qquad \gcd(S,V)=1,
\]
where $S$ is squarefree. All ordinary zeros of $\overline U$
are double, and all its zeros on $D$ are simple.
\end{proposition}
\begin{proof}
Lemma~\ref{lem:multiplicities} determines the multiplicities over
$\overline\Fp$. Zero is not a root, because $\overline U(0)=1$.
Let $g=\gcd(\overline U,\overline U')$ be monic in $\Fp[t]$.
Each double root contributes exactly once to $g$, and a simple
root contributes nothing. Thus $g(0)\ne0$ and
$V=g/g(0)$ lies in $\Fp[t]$ with constant term $1$.
The quotient $S=\overline U/V^2$ is a polynomial with constant
term $1$, whose roots are precisely the remaining simple roots.
They are roots of $D$, so $S\mid D$. They are disjoint from the
roots of $V$, proving the last assertions.
\end{proof}

For the sequence $Q_n$, $D=(1+32t)(1+36t)$ is separable for $p>3$.
An explicit form of $S$ uses, in addition, Sun's classical
\emph{Theorem}~3.5 modulo $p$ \cite[p.~2739]{sun}, which determines
the vanishing of $G_p(-1/36)$ for $p>5$.
It is an auxiliary source input for the next corollary, not a
premise of Theorem~\ref{thm:main}. It is distinct from Sun's
\emph{Conjecture}~3.5 modulo $p^2$.

\begin{corollary}\label{cor:explicit-squarefree}
For every prime $p>3$, there is $V_p\in\Fp[t]$, $V_p(0)=1$, such that
\begin{equation}\label{eq:explicit-squarefree}
 G_p(t)\equiv
 (1+32t)^{e_2(p)}(1+36t)^{e_3(p)}V_p(t)^2\pmod p,
\end{equation}
where
\[
 e_2(p)=\frac{1-\leg{-2}p}{2},\qquad
 e_3(p)=\frac{1-\leg{-3}p}{2}.
\]
\end{corollary}
\begin{proof}
For $p>5$, Sun's Theorem~3.7 determines whether $-1/32$ is a
zero, and Theorem~3.5 determines whether $-1/36$ is a zero.
In each case the inert residue is zero and the split residue is
$4x^2$, nonzero modulo $p$, with its corresponding quadratic
representation. Thus these zeros occur exactly when
$\leg{-2}p=-1$ and $\leg{-3}p=-1$, respectively.
Proposition~\ref{prop:squarefree} gives the factorization and its
normalization at zero. For $p=5$, \eqref{eq:p5-polynomial} equals
$D(t)$ modulo $5$, and both characters are $-1$; take $V_5=1$.
\end{proof}

\subsection{An exact Fourier trace and two Hecke coefficients}
\begin{proposition}\label{prop:Fourier-trace}
The Fourier expansion of the certified period is
\begin{equation}\label{eq:Fourier}
 \Theta(q)=\sum_{n\ge0}c_nq^n,\qquad c_0=1,\qquad
 c_n=-12\sum_{\substack{d\mid n\\(d,6)=1}}d\quad(n\ge1).
\end{equation}
For every prime $p>3$,
\begin{equation}\label{eq:Fourier-trace}
 p\Theta(q^p)+\frac1p\sum_{k=0}^{p-1}
 \Theta(\zeta_p^kq^{1/p})=(p+1)\Theta(q).
\end{equation}
For the polynomial normalization in \eqref{eq:Hecke-product},
\begin{equation}\label{eq:C0-C1}
 C_0=p^{2p},\qquad C_1=-p^{2p-1}(p+1).
\end{equation}
\end{proposition}
\begin{proof}
Logarithmic differentiation of the eta product $R$ in
\eqref{eq:logR} gives
\[
 c_n=-12\bigl(\sigma_1(n)-2\sigma_1(n/2)
                  -3\sigma_1(n/3)+6\sigma_1(n/6)\bigr),
\]
with $\sigma_1$ set to zero on nonintegral arguments.
Writing $n=2^a3^bm$, $(m,6)=1$, and using
$\sigma_1(2^a)-2\sigma_1(2^{a-1})=1$ and the analogous identity
for $3$, the bracket equals $\sigma_1(m)$.
This proves \eqref{eq:Fourier} for every coefficient.

The root-of-unity filter makes the coefficient of $q^n$ on the
left of \eqref{eq:Fourier-trace} equal to
$p c_{n/p}+c_{pn}$, with $c_{n/p}=0$ for $p\nmid n$.
For $n>0$, the geometric sum in the $p$-part of the divisor sum
gives $p c_{n/p}+c_{pn}=(p+1)c_n$.
For $n=0$, both sides are $p+1$.
Finally, the product defining $\Tp$ gives $C_0=p^{2p}$ and
\[
 C_1=-p^{2p-2}
 \frac{p^2\Theta(q^p)+\sum_k\Theta(\zeta_p^kq^{1/p})}{\Theta(q)}.
\]
The trace identity evaluates this as the constant in
\eqref{eq:C0-C1}.
\end{proof}

This trace gives a second method of generating the period coefficients
and checks the normalization of the Hecke product.
The degree bound and the finite specialization still require the
arguments of Sections~\ref{sec:Hecke} and~\ref{sec:finite}.

\subsection{The modulus cannot be raised uniformly to a cube}
The corresponding equality and inert vanishing modulo $p^3$ fail
already at $p=7$. Direct calculation from \eqref{eq:sequences} gives
\begin{equation}\label{eq:p3-limit}
 G_7(-1/32)\equiv294,\qquad G_7(1/64)\equiv49\pmod{343}.
\end{equation}
Both are zero modulo $49$, as Theorem~\ref{thm:main} asserts;
neither is zero modulo $343$, and the two residues are distinct.
This example fixes the scope of the stated supercongruence.

\appendix
\section{Finite certificates for the modular identities}
\label{app:certificates}
The certificate in Lemma~\ref{lem:weight8} is an exact coefficient
calculation coupled with the holomorphic valence bound.
Let $\ell_J=\dd\log J$ and let $A_8,B_8$ denote the left sides
of \eqref{eq:P1} and \eqref{eq:P2}. The necessary coefficients are
\begin{center}\small
\begin{tabular}{r r r r r r r}
\toprule
$n$ & $[q^n]h$ & $[q^n]\Theta$ & $[q^n]J$ & $[q^n]\ell_J$
 & $[q^n]A_8$ & $[q^n]B_8$\\
\midrule
0 & 0 & 1 & 0 & 1 & 0 & 0\\
1 & 1 & $-12$ & 1 & $-2$ & 0 & 0\\
2 & 5 & $-12$ & $-2$ & $-10$ & 0 & 0\\
3 & 19 & $-12$ & $-3$ & $-14$ & 0 & 0\\
4 & 61 & $-12$ & 4 & $-26$ & 0 & 0\\
5 & 174 & $-72$ & 6 & $-12$ & 0 & 0\\
6 & 455 & $-12$ & 6 & $-70$ & 0 & 0\\
7 & 1112 & $-96$ & $-16$ & $-16$ & 0 & 0\\
8 & 2573 & $-12$ & $-8$ & $-58$ & 0 & 0\\
\bottomrule
\end{tabular}
\end{center}
The weight-two comparison in Lemma~\ref{lem:eta-identities} has
\begin{center}
\begin{tabular}{r r r r}
\toprule
$n$ & $[q^n]\Theta$ & $[q^n](-\dd\log R)$
 & $[q^n](\Theta+\dd\log R)$\\
\midrule
0 & 1 & 1 & 0\\
1 & $-12$ & $-12$ & 0\\
2 & $-12$ & $-12$ & 0\\
\bottomrule
\end{tabular}
\end{center}
These tables can be recomputed entirely from the eta products.
For an exponent vector $(r_d)$ write
\begin{equation}\label{eq:eta-product-algorithm}
 \prod_{d\mid6}\eta_d^{r_d}
 =q^{\sum d r_d/24}
  \prod_{d\mid6}\prod_{m\ge1}(1-q^{dm})^{r_d}.
\end{equation}
After accounting for the initial shift, only factors with
$dm\le K$ can affect coefficients through degree $K$.
Expand each factor by the integral binomial series and multiply
with truncation after every product. This is a finite exact algorithm.

Alternatively, put $A(q)=\prod_{d\mid6}\prod_{m\ge1}(1-q^{dm})^{r_d}$,
$A(0)=1$. Its logarithmic derivative is
\[
 \frac{\dd A}{A}=\sum_{n\ge1}\ell_nq^n,
 \qquad
 \ell_n=-\sum_{\substack{d\mid6\\d\mid n}}
              r_d d\,\sigma_1(n/d),
\]
so its coefficients obey the exact recursion
\begin{equation}\label{eq:eta-recurrence-algorithm}
 n A_n=\sum_{j=1}^n\ell_j A_{n-j}.
\end{equation}
For the logarithmic derivative required at degree eight, use directly
\begin{equation}\label{eq:ell-J}
 \ell_J=1-2\sum_{n\ge1}
   \left(\sum_{\substack{d\mid6\\d\mid n}}d\,\sigma_1(n/d)\right)q^n.
\end{equation}
In particular, the coefficient of degree eight of $\ell_J$ is not
inferred from the displayed prefix of $J$ alone.
Ordinary sums, products and the rule $\dd q^n=nq^n$ now reproduce
all residual entries. The table supplies the complete eighteen
coefficients needed at weight eight, while the three weight-two
coefficients suffice at its smaller bound.

\section{Exact eta phases at the two finite cusps}\label{app:cusps}
We derive the signs in \eqref{eq:cusp-values} from the Dedekind
eta law \cite[23.18.5--23.18.7]{dlmf}.
For $c=2,3$, let $\sigma_c=\smat10c1$.
For $d=1,2,3,6$, put $g=\gcd(c,d)$, $a=d/g$, $c'=c/g$,
and choose $e,b\in\Z$ with $ae-bc'=1$.
The exact matrix decomposition is
\begin{equation}\label{eq:cusp-factorization}
 \begin{pmatrix}d&0\\c&1\end{pmatrix}
 =\begin{pmatrix}a&b\\c'&e\end{pmatrix}
  \begin{pmatrix}g&-b\\0&a\end{pmatrix}.
\end{equation}
Let $s(u,v)$ be the Dedekind sum
\[
 s(u,v)=\sum_{j=1}^{v-1}
  \left(\frac jv-\frac12\right)
  \left(\left\{\frac{uj}{v}\right\}-\frac12\right)
 \quad(\gcd(u,v)=1,\ v>0),
\]
where the empty sum for $v=1$ is zero.
All summands here have nonintegral $uj/v$.
Along $\tau=iY$, $Y\to+\infty$, the eta law and its leading term give
\begin{equation}\label{eq:cusp-eta-asymptotic}
 \eta(d\sigma_c\tau)\sim
 \sqrt{\frac gd}\,\sqrt{-i(c\tau+1)}\,
 e^{\pi i\lambda_{c,d}}q^{g^2/(24d)},\qquad
 \lambda_{c,d}=\frac{a+e}{12c'}+s(-e,c')-\frac{b}{12a}.
\end{equation}
The roots are principal. The scale $g/d$ is positive, and
$-i(c\tau+1)$ has positive real part along the stated path,
so the factorization introduces no extra sign.

For $h$ the exponents are $(r_1,r_2,r_3,r_6)=(-5,1,-1,5)$.
The common square-root factor cancels because $\sum r_d=0$,
and the $q$ exponent cancels because both cusp orders are zero.
The following choices give all eight phases explicitly:
\begin{center}\small
\begin{tabular}{r r r r r r r r}
\toprule
$c$ & $d$ & $g$ & $a$ & $c'$ & $e$ & $b$ & $\lambda_{c,d}$\\
\midrule
2 & 1 & 1 & 1 & 2 & 1 & 0 & $1/12$\\
2 & 2 & 2 & 1 & 1 & 1 & 0 & $1/6$\\
2 & 3 & 1 & 3 & 2 & 1 & 1 & $5/36$\\
2 & 6 & 2 & 3 & 1 & 1 & 2 & $5/18$\\
3 & 1 & 1 & 1 & 3 & 1 & 0 & $0$\\
3 & 2 & 1 & 2 & 3 & 2 & 1 & $1/8$\\
3 & 3 & 3 & 1 & 1 & 1 & 0 & $1/6$\\
3 & 6 & 3 & 2 & 1 & 1 & 1 & $5/24$\\
\bottomrule
\end{tabular}
\end{center}
Consequently
\begin{center}
\begin{tabular}{c c c}
\toprule
$c$ & $\sum_{d\mid6}r_d\lambda_{c,d}$
 & $\prod_{d\mid6}(\gcd(c,d)/d)^{r_d/2}$\\
\midrule
2 & 1 & $1/9$\\
3 & 1 & $1/8$\\
\bottomrule
\end{tabular}
\end{center}
In each case the total phase is $e^{\pi i}=-1$.
Therefore $h(1/2)=-1/9$ and $h(1/3)=-1/8$, including their signs.

\section{Companion computations}\label{app:companion}
The companion software, version~0.1.0, implements the original
binomial definitions, finite invariants, local lift checks, eta
products and modular residual calculations, and exact Hecke
polynomial reconstruction. It also provides bounded regression
checks of the sums and the complex multiplication matrices.
Companion repository:
\begin{center}\small
\url{https://github.com/aconsciousfractal/Sun-3-6-Supercongruence}
\end{center}
The repository README and \texttt{REPRODUCE.md} give the verification
and build instructions. The coefficient algorithms in
Appendix~\ref{app:certificates} and the exact phase tables in
Appendix~\ref{app:cusps} specify the finite certificates used in
the proof. Tests over bounded sets of primes serve as regression
checks; the all-prime conclusion follows from the arguments above.
The evidence inventory is \texttt{EVIDENCE\_SHA256.txt}; its SHA-256 is
\begin{center}\footnotesize\ttfamily
\detokenize{47f214dfbf8133220e57c1a67ade26d09c0343fe17df719e38ad98c26a86f1c8}
\end{center}
The inventory identifies companion code, tests, configuration and
exact evidence, excluding the manuscript and prose that quote the digest.

\begin{thebibliography}{9}
\bibitem{sun}
Zhi-Hong Sun,
\emph{Supercongruences involving Ap\'ery-like numbers and binomial coefficients},
AIMS Mathematics \textbf{7}(2) (2022), 2729--2781.
\href{https://doi.org/10.3934/math.2022153}{doi:10.3934/math.2022153}.
Definitions and recurrence, pp.~2729--2730;
Theorems~3.5 and~3.7, pp.~2739--2740;
Conjecture~3.6, p.~2742.

\bibitem{beukers}
Frits Beukers,
\emph{Supercongruences using modular forms},
arXiv:2403.03301v3, 16 June 2025.
\url{https://arxiv.org/abs/2403.03301v3}.
Section~3, pp.~16--18, especially Lemma~3.7;
Section~5, pp.~20--22; Theorem~6.1, p.~24; Section~8.10, p.~28.

\bibitem{straub}
Armin Straub,
\emph{Gessel--Lucas congruences for sporadic sequences},
arXiv:2301.12248v1, 28 January 2023.
\url{https://arxiv.org/abs/2301.12248v1}.
Equation~(10) and Theorem~3.2, p.~5; proof, pp.~6--7.
The positive sequence is $A_F(n)=(-1)^nQ_n$.
Journal publication: Monatshefte f\"ur Mathematik \textbf{203}
(2024), 883--898,
\href{https://doi.org/10.1007/s00605-023-01894-3}{doi:10.1007/s00605-023-01894-3}.
The proof locators refer to the specified arXiv version.

\bibitem{gorodetsky}
Ofir Gorodetsky,
\emph{New representations for all sporadic Ap\'ery-like sequences,
with applications to congruences},
arXiv:2102.11839v2, 5 January 2025.
\url{https://arxiv.org/abs/2102.11839v2}.
Proposition~1.2, equation~(1.8), p.~4;
Proposition~3.3 and row F, p.~10; multinomial derivation, p.~12.
Journal publication: Experimental Mathematics \textbf{32}(4)
(2023), 641--656,
\href{https://doi.org/10.1080/10586458.2021.1982080}{doi:10.1080/10586458.2021.1982080}.
The locators above refer to the specified arXiv version.

\bibitem{rousewebb}
Jeremy Rouse and John J. Webb,
\emph{On spaces of modular forms spanned by eta-quotients},
arXiv:1311.1460v3, 14 November 2014.
\url{https://arxiv.org/abs/1311.1460v3}.
Eta-quotient criterion and cusp-order formula, p.~2.
Journal publication: Advances in Mathematics \textbf{272} (2015), 200--224.
\href{https://doi.org/10.1016/j.aim.2014.12.002}{doi:10.1016/j.aim.2014.12.002}.

\bibitem{craig}
William Craig,
\emph{New Types of Sturm bounds via $p$-adic transfer methods},
arXiv:2602.10240v1, 10 February 2026.
\url{https://arxiv.org/abs/2602.10240v1}.
Only the classical holomorphic-form bound in
Theorem~1.1(1), p.~2, is used.

\bibitem{dlmf}
NIST Digital Library of Mathematical Functions,
Section~23.18, \emph{Modular Transformations},
equations 23.18.5--23.18.7.
\url{https://dlmf.nist.gov/23.18}.

\bibitem{chancooper}
Heng Huat Chan and Shaun Cooper,
\emph{Rational analogues of Ramanujan's series for $1/\pi$},
Mathematical Proceedings of the Cambridge Philosophical Society
\textbf{153}(2) (2012), 361--383.
\href{https://doi.org/10.1017/S0305004112000254}{doi:10.1017/S0305004112000254}.
The level-$6$ parametrization is cited for comparison; it is not
needed for the direct period certification.
\end{thebibliography}
\end{document}
